\documentclass[11pt]{article}

\usepackage[a4paper,margin=1in]{geometry}
\usepackage[T1]{fontenc}
\usepackage[utf8]{inputenc}
\usepackage{lmodern}
\usepackage{amsmath,amssymb}

\title{Saturated Superfluid Vacuum III\\
Irreversible Time and Wake Entropy}
\author{Stig Norland}
\date{Draft}

\begin{document}

\maketitle

\begin{abstract}
This paper addresses only the origin of irreversible time in the Saturated Superfluid Vacuum
(SSV) framework. We argue that time is not a background dimension through which reality
moves, but the local rate of change of the vacuum order parameter $\Psi$. The past does not
exist as a destination or region; it survives only as wake structure encoded in the present
state of the medium. Events in $\Psi$ do not simply occur and vanish. They deposit
distributed traces in the form of phonon phases, vortex distortions, Kelvin-wave content,
density correlations, and altered future reconnection opportunities. We therefore propose
that entropy in SSV should not be defined as a bare event count, but as the persistent
wake-complexity of the present state.

On this view, irreversible time is not the flow of an external parameter, but the one-way
accumulation and redistribution of wake structure in $\Psi$. Formal time reversal in equations
is distinguished from physical reversal in the medium: the substitution $t \to -t$ defines a
mathematical transformation on trajectories, but does not supply a physical operation that
globally unwrites the distributed wake of an event. This dissolves the usual paradoxes of time
travel by treating them as category errors arising from the reification of the past as a place.
\end{abstract}

\newpage
\tableofcontents
\newpage

\section{Introduction}

The standard treatment of time in fundamental physics is powerful, but conceptually unstable.
At the level of formal equations, time often appears as a reversible parameter. At the level of
experience and thermodynamics, it appears as an irreversible order of change. This tension is
usually managed rather than resolved.

The aim of this paper is narrow. We address only the origin of irreversible time in the SSV
framework introduced in Paper~I~\cite{SSV-I} and developed in Paper~II~\cite{SSV-II}. The
central claim is:

\begin{quote}
Time is the local rate of change of $\Psi$, the past survives only as wake structure in the
present state, and entropy is the persistent wake-complexity of that state.
\end{quote}

This paper is therefore about ontology first and thermodynamic interpretation second. The
goal is to state the time claim cleanly enough that the later papers can inherit it without
having to carry the whole argument again.

\subsection{Roadmap}

Section 2 states the minimal SSV setting used in this paper. Section 3 states the ontological
claim that time is local change rather than a background dimension, and draws the immediate
consequence for time-travel paradoxes. Section 4 develops the idea that events write wake
structure into the medium. Section 5 reframes entropy as wake-complexity. Section 6
distinguishes formal time reversal in equations from physical reversal in the medium.
Section 7 closes with the remaining open technical questions.

\section{Minimal SSV Setting}

Only a small part of the SSV framework is needed here. The physical vacuum is represented by a
macroscopic order parameter
\begin{equation}
  \Psi(\mathbf{x},t)=\sqrt{\rho(\mathbf{x},t)}\,e^{i\theta(\mathbf{x},t)},
\end{equation}
with density $\rho$, phase $\theta$, and superfluid velocity proportional to $\nabla\theta$ away
from vortex cores. The exact field equation is not required for the present argument. What is
required is that the state of the medium has real local structure: phase, density, vortex
geometry, phonon content, correlations, and boundary data.

In this paper, an ``event'' means a localized update of that structure whose forward
realization deposits wake content the medium cannot locally re-gather into the prior state.
This criterion is deliberately structural rather than threshold-based. It admits any process
whose forward and reverse boundary-data requirements differ — a reconnection, a phonon
emission, a Kelvin-wave excitation, a vortex displacement that radiates — and excludes
purely reversible relabellings and stationary zero-point fluctuations, whose statistics by
construction contribute no net wake to $\mathcal{P}(t)$. Reconnection is the clearest
illustrative case because it changes the connectivity of vortex filaments and converts
localized order into distributed excitations, but it is not the only possible event, and it is
not used here as a primitive counter for entropy.

The criterion is circular in form — wake is what events deposit, events are what deposit wake
— but the circularity is shared with any account of irreversibility, and it is closed
empirically by the candidate functional in \S5: a process counts as an event to the extent
that it raises $W[\Psi]$ in any of its channels.

The present state of the medium will be denoted schematically by
\begin{equation}
  \mathcal{P}(t)=\{\rho,\theta,\mathcal{V},\mathcal{K},\mathcal{C},\ldots\}_t,
\end{equation}
where $\mathcal{V}$ denotes vortex geometry, $\mathcal{K}$ Kelvin-wave content, and
$\mathcal{C}$ correlation data. This notation is deliberately schematic. Its role is to keep the
argument attached to present physical structure rather than to an abstract external time axis.

\section{Time as Change, Not Dimension}

In SSV, the physical vacuum is a real medium described by an order parameter $\Psi$. What
physics ordinarily measures as time is not an independent dimension in which the medium is
embedded, but the local rate at which the state of the medium changes. This is closer in
spirit to the relational treatments of time advanced by Barbour~\cite{Barbour1999} and
Rovelli~\cite{Rovelli1993,Rovelli2018} than to the block-universe reading common in
spacetime physics.

A useful schematic statement is
\begin{equation}
  d\tau \propto \mathcal{R}[\Psi]\,dt,
\end{equation}
where $\mathcal{R}[\Psi]$ denotes the local update-rate of the medium. In later papers this
quantity can be connected to gravitational and kinematic slowing. Here it is used only to
express the ontological point: proper time is not a substance added to $\Psi$, but a measure of
how much local physical change occurs in the medium.

This is not yet a final derivation, but a statement of ontology. The present is the current
state of $\Psi$. The future does not yet exist. The past is not a region one could return to. It
exists only insofar as earlier events have left structure in the present state of the medium.

\subsection{Present-State Ontology and Causal Ordering}\label{sec:presentstate}

The strongest version of the claim has three parts:
\begin{enumerate}
  \item only the present state of $\Psi$ physically exists,
  \item the past is encoded in that state as record, wake, and correlation,
  \item there is no second arena in which earlier states continue to exist as accessible places.
\end{enumerate}

This does not deny ordering. Earlier and later remain meaningful because wake structure is
asymmetric: a present configuration can contain records of a previous event (outgoing phases,
displaced vortex geometry, thermalized phonons, modified reconnection opportunities), while the
pre-event configuration does not contain the later wake as an already-existing region. Causal
order is carried by structure in $\mathcal{P}(t)$, not by an external archive of completed
moments.

A direct consequence is that time-travel paradoxes are not exotic physical possibilities but
signs of a mistaken description. They rely on the assumption that the past remains available as
a destination; in SSV there is no such region, and a journey ``to'' the past would require the
current state of $\Psi$ to be both itself and an earlier state whose records it already
contains. That is a contradiction in description, not a physical process to be analyzed.

\section{Wake Production in $\Psi$}

An event in $\Psi$ does not simply happen and disappear. It perturbs the medium and leaves a
wake. Depending on scale and context, that wake may include:

\begin{itemize}
  \item outgoing phonon phases,
  \item Kelvin-wave excitations,
  \item vortex displacement and writhe,
  \item density ripples,
  \item altered local reconnection possibilities,
  \item longer-range correlations in the surrounding medium.
\end{itemize}

The important point is not that every trace remains separately visible forever. The point is
that the event is redistributed into the present state of the medium rather than physically
unwritten.

\subsection{The Wake Is Distributed}

The wake of an event is generally not stored in one location. A reconnection, for example, may
leave a local change in vortex geometry, a packet of outgoing phonons, a spectrum of Kelvin
waves on the retracted filaments, and altered phase correlations in the surrounding medium.
After further interactions these traces may become individually hard to identify, but their
information has not thereby been gathered into the unique configuration needed to restore the
pre-event state.

This distributed character is what makes wake entropy different from a memory metaphor. The
medium does not merely ``remember'' in an observer-dependent sense. Its present physical state
has been changed in ways that affect what subsequent local updates are possible.

\subsection{Why Reconnection Matters But Does Not Exhaust the Story}\label{sec:reconnection}

Vortex reconnection remains an important process in SSV because it is one of the clearest
mechanisms by which localized structure is converted into distributed wake structure. A
reconnection event changes topology, emits phonons, populates Kelvin-wave sectors, and alters
future reconnection possibilities. But the present paper does not identify entropy with a bare
count of reconnections. Reconnection is one important wake-writing process inside a broader
picture in which any event that leaves persistent structure in $\Psi$ contributes to the
irreversible burden carried by the present state.

\subsection{Local Cancellation Versus Global Unwriting}

Local amplitudes can interfere destructively. Some observables can relax, thermalize, or become
coarse-grained away. But this does not by itself restore the pre-event global state. To reverse
an event physically would require the distributed wake to be gathered and re-focused into the
unique incoming configuration that recreates the earlier state.

That is a nonlocal requirement. No local dynamical process in the medium is known to perform
it.

\section{Entropy as Wake-Complexity}

This motivates a reframing of entropy in SSV. Entropy should not be introduced as a bare count
of past events. That would risk tautology if monotonicity is built into the definition from the
start. The stronger claim is:

\begin{quote}
Entropy is the irreducible wake-complexity of the present state of $\Psi$.
\end{quote}

\subsection{Why This Is Not Tautological}\label{sec:nottautological}

The monotonicity claim is not:
\begin{quote}
Entropy increases because we defined it to count only positive increments.
\end{quote}

It is instead:
\begin{quote}
Events write distributed structure into $\Psi$, and no physical process available to the medium
globally unwrites that structure back into the exact pre-event state.
\end{quote}

This is a physical claim about the availability of inverse processes, not merely a bookkeeping
choice.

\subsection{Relation to Boltzmann Entropy}\label{sec:boltzmann}

The standard expression $S=k_B\ln\Omega$ remains useful; SSV changes the interpretation of the
microstates being counted. A statistical form for wake entropy is
\begin{equation}
  S_{\text{wake}}(t) = k_B \ln \Omega_{\text{wake}}[\Psi(t)],
\end{equation}
where $\Omega_{\text{wake}}$ counts physically distinct present wake configurations compatible
with the same coarse description, rather than an abstract set of hidden possibilities beneath a
macrostate.

The phrase ``the same coarse description'' has to do real work. Concretely, ``the same''
means: equivalent under the smoothing of $\Psi$ to some scale $\ell$ and under restriction to
a specified observable set $\mathcal{O}$. Two wake configurations are identified when their
$\Psi$-fields, smoothed at scale $\ell$, agree on every observable in $\mathcal{O}$. Different
choices of $(\ell,\mathcal{O})$ generate different $\Omega_{\text{wake}}$ and therefore
different numerical entropies — a familiar feature of coarse-grained entropy, made explicit
here so that $\Omega_{\text{wake}}$ is not a rename of $\Omega$ but a counting tied to the
medium's resolvable wake channels.

This also explains why entropy can increase while the microscopic state remains definite. The
medium has one actual state, but the macroscopic description loses access to the detailed
distribution of wake structure. Thermodynamic entropy is then the coarse-grained shadow of a
real structural burden in $\Psi$.

\subsection{Candidate Wake Functional}\label{sec:candidatefunctional}

Equivalently, one may write
\begin{equation}
  S_{\text{wake}}(t) = k_B\, W[\Psi(t)],
\end{equation}
where $W[\Psi(t)]$ measures the amount of present-state structure required to encode the
surviving wake of prior events. A usable functional should respond to the parts of $\Psi$ that
carry persistent wake structure. One possible schematic decomposition is
\begin{equation}
  W[\Psi] =
  \alpha_{\phi} I_{\phi}
  + \alpha_{\rho} I_{\rho}
  + \alpha_{\mathcal{V}} I_{\mathcal{V}}
  + \alpha_{\mathcal{K}} I_{\mathcal{K}}
  + \alpha_{\mathcal{C}} I_{\mathcal{C}},
\end{equation}
where the $I$ terms measure phase disorder, density disorder, vortex tangle complexity,
Kelvin-wave excitation, and correlation complexity respectively. The coefficients $\alpha_i$
set the relative thermodynamic weight of each channel.

Any acceptable $I_i$ must satisfy three requirements. It must be a functional of the present
state, it must be insensitive to purely reversible relabellings, and it must increase when
localized order is converted into distributed wake structure that cannot be locally
unwritten.

This subsection states the decomposition as a checklist; the five subsections
\S\ref{sec:phononreduction}--\S\ref{sec:densityreduction} then carry it out, defining each
$I_i$ as a concrete dimensionless functional, fixing each coefficient $\alpha_i$ by
calibration to a known thermodynamic or counting limit, and reducing each channel to a
recognizable entropy. The outcome is summarized in \S\ref{sec:channelcomplete}. The reader
who wants the schematic picture first may treat the present equation as a checklist and read
those subsections as its execution.

\subsection{Fate of the Wake Channels}\label{sec:channelfate}

Although the $I_i$ are not yet canonical functionals, their physical content is concrete
enough to classify by fate. Each channel of $W[\Psi]$ corresponds to wake content with a
characteristic mode of propagation, decay, or persistence in the medium, and the irreversible
burden carried by the present state is not distributed evenly across them. Table~1 summarizes
the dominant fate of each channel.

\begin{table}[h]
\centering
\begin{tabular}{lll}
\hline
Channel & Dominant fate & Notes \\
\hline
$I_\phi$ (phase disorder) & radiated, then thermalized & Carried outward by phonons; \\
 & & mixed into phonon bath on scattering timescales. \\[2pt]
$I_\rho$ (density disorder) & radiated, then thermalized & Sound-wave content; integral mass \\
 & & conservation gives a global constraint only. \\[2pt]
$I_\mathcal{V}$ (vortex tangle) & conserved between events & Line length, writhe, and topology \\
 & & approximately conserved between reconnections; \\
 & & converted at events into $I_\mathcal{K}$ and $I_\phi$. \\[2pt]
$I_\mathcal{K}$ (Kelvin-wave) & radiated along filaments & Kelvin cascade transports to small scales; \\
 & & ultimately decays into $I_\phi$ via Kelvin-phonon \\
 & & coupling at short wavelength. \\[2pt]
$I_\mathcal{C}$ (correlation) & erased from observation & Long-range entanglement with the \\
 & & environment; structurally present but typically \\
 & & below resolved $(\ell,\mathcal{O})$. \\
\hline
\end{tabular}
\caption{Dominant fate of each wake channel under forward evolution in $\Psi$. ``Radiated''
means the content propagates outward but remains within $\Psi$; ``thermalized'' means it has
mixed with bath modes such that the macroscopic description loses access to it; ``erased from
observation'' means it persists in $\Psi$ but falls below the coarse-graining
$(\ell,\mathcal{O})$ of any realistic measurement.}
\end{table}

Three structural points follow. First, no channel is genuinely \emph{lost} — every fate
column other than ``conserved'' transfers wake content to another channel or to a sector
below resolution, never out of $\Psi$. Second, vortex tangle $I_\mathcal{V}$ is the only
channel that carries wake across long timescales without converting; this is why reconnection
acts as the rate-limiting event for irreversibility in vortex-dominated regimes (\S\ref{sec:reconnection}). Third, the
classification is observer-relative only in the last column: the wake content tagged as
``erased from observation'' is real structure in $\Psi$ that is missed by a specific
$(\ell,\mathcal{O})$, not absent from the medium. The first quantity of
\S\ref{sec:twoquantities} retains it; the second loses it.

\subsection{Extensivity, Coarse-Graining, and Present Structure}\label{sec:extensivity}

The wake formulation is not meant to deny the usefulness of ordinary thermodynamic entropy. It
proposes a deeper physical interpretation for it: standard thermodynamic entropy is recovered
at the macroscopic level by coarse-graining over the fine wake structure of the medium. On this
view, extensivity and additivity are not primitive axioms, but emergent properties of how wake
structure composes across weakly coupled subsystems.

\subsection{Worked Example: Phonon-Bath Reduction}\label{sec:phononreduction}

The preceding claim is concrete enough to demonstrate on a single channel. Consider a region
of $\Psi$ supporting a thermal bath of Bogoliubov phonons at temperature $T$. The Madelung
phase $\theta(\mathbf{x},t)$ decomposes into mode amplitudes; each long-wavelength mode is a
harmonic oscillator with frequency $\omega_k = c_s k$ and thermal occupation
\begin{equation}
  n_k = \frac{1}{e^{\hbar\omega_k/k_B T} - 1}.
\end{equation}

For the phonon channel, take $I_\phi$ to count the per-mode von Neumann entropy of the
bulk acoustic excitations,
\begin{equation}
  I_\phi \;=\; \sum_k \left[(1+n_k)\ln(1+n_k) - n_k \ln n_k\right].
\end{equation}
This is dimensionless, manifestly a functional of the present state of $\Psi$, and vanishes
on the zero-temperature ground state ($n_k = 0$ for all $k$). Note that for Bogoliubov
phonons, the phase and density quadratures are conjugate quadratures of the \emph{same}
oscillator, so a separate ``density'' channel $I_\rho$ would double-count on this sector;
$I_\rho$ should therefore be reserved for non-phonon density structure (defects, solitons,
condensate-density inhomogeneity) outside the bulk acoustic bath.

Setting $\alpha_\phi = k_B$ and applying the coarse-graining $(\ell,\mathcal{O})$ that
resolves mode occupations but discards intra-mode phase information,
\begin{equation}
  \alpha_\phi I_\phi
  \;=\; k_B \sum_k \left[(1+n_k)\ln(1+n_k) - n_k \ln n_k\right]
  \;=\; S_{\text{thermal}},
\end{equation}
which is the standard thermodynamic entropy of a Bose gas of phonons. The familiar
$S \propto T^3$ scaling in the acoustic limit follows from $D(\omega) \propto \omega^2$ in
three dimensions; no further input is required.

Two structural lessons. First, the calibration $\alpha_\phi = k_B$ is fixed by matching to a
known thermodynamic entropy, illustrating in concrete terms the normalization route promised
in \S\ref{sec:candidatefunctional}. Second, the reduction is exact at the level of the phonon
channel and serves as a
template: an analogous construction in each remaining channel (with the appropriate basis for
$I_\mathcal{K}$, $I_\mathcal{V}$, $I_\mathcal{C}$, and a non-phonon $I_\rho$) would reduce
$W[\Psi]$ to the full thermodynamic entropy under the union of those coarse-grainings. The
next two subsections carry out the reduction on two further independent channels, the Kelvin
and vortex-tangle sectors.

\subsection{Worked Example: Kelvin Sector}\label{sec:kelvinreduction}

The Kelvin sector is genuinely independent of the bulk phonon bath: Kelvin waves are
transverse bending excitations of quantized vortex filaments, supported on the one-dimensional
locus of the vortex core rather than in the three-dimensional bulk. Their dispersion in the
long-wavelength limit is~\cite{Vinen2002}
\begin{equation}
  \omega_\mathcal{K}(k) \;\approx\; \frac{\hbar}{2 m}\, k^2 \,\Lambda(k\xi),
  \qquad \Lambda(k\xi) \;=\; \ln(1/k\xi) + O(1),
\end{equation}
with $m$ the boson mass and $\xi$ the healing length; the slow logarithm reflects the
long-range interaction along the filament. In thermal equilibrium each Kelvin mode is again
a harmonic oscillator with occupation
\begin{equation}
  n_k^\mathcal{K} \;=\; \frac{1}{e^{\hbar \omega_\mathcal{K}(k)/k_B T} - 1}.
\end{equation}

Take $I_\mathcal{K}$ to count the per-mode von Neumann entropy of the Kelvin sector,
\begin{equation}
  I_\mathcal{K} \;=\; \sum_k
  \left[(1 + n_k^\mathcal{K})\ln(1 + n_k^\mathcal{K}) - n_k^\mathcal{K} \ln n_k^\mathcal{K}\right],
\end{equation}
where the sum runs over Kelvin modes on the filament. Setting $\alpha_\mathcal{K} = k_B$ and
applying the analogous coarse-graining $(\ell, \mathcal{O})$ that resolves Kelvin-mode
occupations but discards intra-mode phases,
\begin{equation}
  \alpha_\mathcal{K} I_\mathcal{K} \;=\; S_{\text{Kelvin}}^{\text{thermal}},
\end{equation}
the standard thermodynamic entropy of a Bose gas of Kelvin quanta on the filament.

Two remarks. First, although the formula for the per-mode entropy is identical to the phonon
case, the underlying mode basis is different: bulk acoustic modes labeled by $\mathbf{k}
\in \mathbb{R}^3$ for $I_\phi$ versus one-dimensional filament-bending modes for
$I_\mathcal{K}$. The sum $\alpha_\phi I_\phi + \alpha_\mathcal{K} I_\mathcal{K}$ is therefore
genuinely additive and does \emph{not} double-count, in contrast to the conjugacy issue
flagged in \S\ref{sec:phononreduction} for a putative $I_\rho$ on the same phonon modes.
Second, because $\omega_\mathcal{K}(k) \propto k^2 \Lambda$ rather than $c_s k$, the
low-temperature scaling of $S_{\text{Kelvin}}^{\text{thermal}}$ differs from the acoustic
$T^3$ law; the Kelvin channel therefore contributes a thermodynamically distinct dependence
on $T$ at low temperature, which is in principle observable in geometries where the Kelvin
sector dominates over bulk phonons (slender vortex filaments at low density).

\subsection{Worked Example: Vortex-Tangle Sector}\label{sec:tanglereduction}

The vortex-tangle channel differs in kind from the two preceding examples. Phonons and Kelvin
waves are harmonic excitations, and their entropy is the thermal entropy of a Bose gas of
modes. A vortex tangle is not harmonic: its wake content is \emph{configurational and
topological} --- the geometric arrangement of the filament network, its total quantized line
length, and the knotting and linking of its components. The appropriate entropy is therefore
a configurational count, and this channel turns out to be the most direct realization of the
$S = k_B \ln\Omega_{\text{wake}}$ form of \S\ref{sec:boltzmann}.

Discretize the medium at the healing length $\xi$, which is the physical core size of a
quantized vortex and the natural short-distance cutoff: two filament segments approaching
within $\sim\xi$ either reconnect or cease to be independent. A tangle of total line length
$L$ then occupies $\ell = L/\xi$ filament segments on a lattice of spacing $\xi$. Because
vortex cores cannot overlap, the filament network is a self-avoiding set of closed loops
(together with any open arcs terminating on boundaries). The number of distinguishable tangle
configurations at fixed total length scales, to leading order, as
\begin{equation}
  \Omega_\mathcal{V} \;\sim\; \mu_{\text{eff}}^{\,L/\xi},
\end{equation}
where $\mu_{\text{eff}}$ is a dimensionless constant of order unity --- the connective
constant of the self-avoiding filament ensemble, enlarged to absorb the exponentially growing
multiplicity of distinct knot and link types accessible at length $L/\xi$. Its precise value
is regularization-dependent (for self-avoiding walks on a simple cubic lattice the geometric
part alone gives $\mu \approx 4.68$); what is robust is that $\ln\mu_{\text{eff}}$ is a
positive number of order unity.

The tangle-channel functional is then
\begin{equation}
  I_\mathcal{V} \;=\; \ln\Omega_\mathcal{V} \;=\; \frac{L}{\xi}\,\ln\mu_{\text{eff}}
  \;+\; (\text{subleading loop-closure and loop-number terms}).
\end{equation}
This is dimensionless ($L/\xi$ is a pure ratio), is a functional of the present state of
$\Psi$ --- the core locus is the set where $\rho \to 0$ with nonzero phase winding ---
vanishes on the vortex-free ground state ($L = 0$), and is invariant under rigid relabellings
of the medium.

Setting $\alpha_\mathcal{V} = k_B$, as for the previous two channels,
\begin{equation}
  \alpha_\mathcal{V} I_\mathcal{V} \;=\; k_B \ln\Omega_\mathcal{V} \;=\; S_{\text{config}},
\end{equation}
a bona fide Boltzmann configurational entropy of the vortex tangle. No coarse-graining over a
mode bath is required here: the tangle configurations \emph{are} the microstates, and the
vortex channel is therefore the cleanest instance of the interpretation advanced in
\S\ref{sec:boltzmann}, where $\Omega_{\text{wake}}$ counts physically distinct present
configurations rather than abstract hidden possibilities beneath a macrostate.

One subtlety is essential and, properly understood, strengthens the framework: \emph{the
tangle functional $I_\mathcal{V}$ is not individually monotonic.} In freely decaying quantum
turbulence the line-length density falls --- in the Vinen regime as $L \propto
t^{-1}$~\cite{Vinen2002} --- so $I_\mathcal{V}$ \emph{decreases} with time. This does not
violate the monotonicity of $W[\Psi]$. Tangle decay proceeds by reconnection, and each
reconnection converts organized line length into Kelvin-wave excitation and radiated phonons:
wake flows out of $I_\mathcal{V}$ and into $I_\mathcal{K}$ and $I_\phi$. The total
$W = \sum_i \alpha_i I_i$ does not decrease, because the conversion is dissipative --- an
organized, low-entropy channel feeding high-entropy thermalized channels. Only the sum is
protected; the individual channels exchange wake. This is exactly the fate-table entry of
\S\ref{sec:channelfate}, in which $I_\mathcal{V}$ carries wake across long timescales until
reconnection converts it, and it is why reconnection is the rate-limiting event for
irreversibility in vortex-dominated regimes.

Finally, $I_\mathcal{V}$ is directly computable. The total line length $L$ is a
straightforward readout of any simulated state of $\Psi$: locate the core locus as the set of
density zeros carrying nonzero phase winding, and measure its length. The knotted-filament
states and continuation sweeps of Paper~I~\cite{SSV-I} are therefore immediate test data for
the vortex channel, and a measured $I_\mathcal{V}(t)$ along a reconnection cascade --- paired
with $I_\mathcal{K}(t)$ and $I_\phi(t)$ --- would exhibit the inter-channel wake transfer
just described as an explicit budget that is conserved in sum.

\subsection{Worked Example: Correlation Channel}\label{sec:correlationreduction}

The four channels treated so far are all \emph{local} in character: each counts entropy
carried by degrees of freedom identified at a point or along a filament --- bulk acoustic
modes, Kelvin modes, vortex line. The correlation channel $I_\mathcal{C}$ is different in
kind. It measures wake that resides in no single region but in the \emph{joint statistics} of
separated regions: the imprint left when one event correlates parts of the medium that a
local description would treat as independent.

The natural measure is mutual information. Partition the coarse-grained medium at scale
$\ell$ into regions, and for two regions $A$ and $B$ at separation $r$ define
\begin{equation}
  I(A\!:\!B) \;=\; S(A) + S(B) - S(A\cup B),
\end{equation}
with $S(\cdot)$ the entropy of the coarse-grained reduced state. A subtlety must be settled
before this can serve as a channel of $W[\Psi]$. A state in local thermal equilibrium
already carries correlations --- the equilibrium correlations dictated by the same mode
occupations already counted in $I_\phi$ and $I_\mathcal{K}$. Including those in
$I_\mathcal{C}$ would double-count. The correlation channel must therefore count only the
\emph{connected, long-range} mutual information: the part that remains once the contributions
already attributed to the local channels are removed, and that persists at separations $r$
much larger than the healing length $\xi$ and the thermal correlation length $\lambda_T$.
Write
\begin{equation}
  I_\mathcal{C} \;=\; \sum_{\text{region pairs}} I_{\text{conn}}(A\!:\!B),
  \qquad I_{\text{conn}}(A\!:\!B)\big|_{r \,\lesssim\, \xi,\,\lambda_T} \;\equiv\; 0,
\end{equation}
the sum running over pairs in the partition family, with the short-range equilibrium part
subtracted by construction.

So defined, $I_\mathcal{C}$ is dimensionless, is a functional of the present state of
$\Psi$, and vanishes on any state of local thermal equilibrium. Setting
$\alpha_\mathcal{C} = k_B$,
\begin{equation}
  \alpha_\mathcal{C} I_\mathcal{C}
  \;=\; k_B \sum_{\text{region pairs}} I_{\text{conn}}(A\!:\!B),
\end{equation}
a genuine entropy in thermodynamic units.

This channel carries the conceptual payoff of the whole decomposition. The first three
worked examples each reduced a channel to a recognizable entropy --- two Bose-gas entropies
and a Boltzmann configurational entropy. The correlation channel does the opposite: it is the
channel that \emph{has no counterpart in ordinary local thermodynamics}, and that is exactly
why it is the formal home of the wake idea. Standard thermodynamic entropy is recovered as
the limit
\begin{equation}
  I_\mathcal{C} \to 0
  \quad\Longrightarrow\quad
  W[\Psi] \to
  \alpha_\phi I_\phi + \alpha_\rho I_\rho
  + \alpha_\mathcal{V} I_\mathcal{V} + \alpha_\mathcal{K} I_\mathcal{K},
\end{equation}
the sum of local channel entropies --- the standard, extensive, locally-additive
thermodynamic entropy of \S\ref{sec:extensivity}. A medium in local equilibrium has
$I_\mathcal{C} = 0$ and obeys ordinary thermodynamics. A medium carrying wake has
$I_\mathcal{C} > 0$: its separated regions share records of a common past. The correlation
channel is the precise statement that the past survives as present structure --- the
structure being the connected correlations that local thermodynamics discards.

Two remarks. First, $I_\mathcal{C}$ is the most observer-relative channel: mutual information
depends on the partition $(\ell,\mathcal{O})$, so $I_\mathcal{C}$ belongs squarely to the
second of the two quantities distinguished in \S\ref{sec:twoquantities} --- the
coarse-grained shadow, not the partition-independent burden carried by $\Psi$ itself.
Second, like $I_\mathcal{V}$, $I_\mathcal{C}$ is not individually monotonic: wake
correlations spread as the medium evolves, raising the number of correlated region pairs
while diluting the mutual information carried by each. Only the total $W$ is protected; the
correlation channel exchanges wake with the others as records disperse.

\subsection{Worked Example: Solitonic Density Channel}\label{sec:densityreduction}

The final channel is the non-phonon density channel $I_\rho$. As noted in
\S\ref{sec:phononreduction}, density and phase are conjugate quadratures of the same
Bogoliubov oscillator, so $I_\rho$ must \emph{not} count small-amplitude acoustic structure
--- that is already in $I_\phi$. What it counts is the genuinely non-perturbative localized
density structure of the medium: solitonic objects --- dark and grey solitons, density
notches, and other codimension-one density defects --- that are neither sound nor vortex
line. They are identified in $\Psi$ by a density depletion exceeding a fixed fraction of
$\bar\rho$, localized over a width of order the healing length $\xi$, and carrying no net
circulation (which would make them vortices, counted in $I_\mathcal{V}$).

The channel is configurational, and the counting follows the template of
\S\ref{sec:tanglereduction}. Consider $N$ such solitonic objects in volume $V$. Each occupies
a phase-space cell $v_0 = c_0\,\xi^{d}$, with $c_0$ a dimensionless constant of order unity
and $d$ the spatial dimension, and carries an internal-parameter multiplicity $g_{\text{int}}$
--- the number of distinguishable depth and velocity states of the soliton family. The number
of distinguishable configurations of the dilute soliton gas is
\begin{equation}
  \Omega_\rho \;=\; \frac{1}{N!}\left(\frac{V}{v_0}\right)^{\!N} g_{\text{int}}^{\,N},
\end{equation}
the $N!$ removing permutations of identical objects. With Stirling's approximation the
channel functional is
\begin{equation}
  I_\rho \;=\; \ln\Omega_\rho
  \;=\; N\left[\ln\!\left(\frac{V}{N v_0}\right) + 1 + \ln g_{\text{int}}\right].
\end{equation}
This is dimensionless, is a functional of the present state of $\Psi$ (soliton number and
parameters are readouts of the density and phase fields), vanishes on the soliton-free ground
state ($N = 0$), and is invariant under rigid relabellings.

Setting $\alpha_\rho = k_B$,
\begin{equation}
  \alpha_\rho I_\rho \;=\; k_B \ln\Omega_\rho \;=\; S_{\text{soliton gas}},
\end{equation}
which is the configurational entropy of a dilute gas of localized objects --- the
Sackur--Tetrode configurational term, with $v_0$ in place of the thermal phase-space cell.
The reduction to a recognizable thermodynamic entropy is therefore immediate, as for the
vortex channel.

Two remarks complete the parallel with the earlier channels. First, $I_\rho$ does not
double-count: solitonic objects are non-perturbative and codimension-one, distinct both from
the linearized phonon bath ($I_\phi$) and from the codimension-two circulating defects of
$I_\mathcal{V}$. Second, $I_\rho$ is not individually monotonic. In three dimensions dark
solitons are unstable to the snake instability and decay into vortex rings --- a wake
transfer from $I_\rho$ into $I_\mathcal{V}$ --- while soliton collisions radiate phonons into
$I_\phi$. As before, $N$ can fall, $I_\rho$ with it, and only the total $W$ is protected.

\subsection{The Channel Set Is Complete}\label{sec:channelcomplete}

With the density channel in hand, all five channels of the candidate functional
(\S\ref{sec:candidatefunctional}) have been given concrete dimensionless definitions and
calibrated coefficients. A structural result emerges. Each $I_i$ was defined as a genuine
entropy measured in nats --- a Bose-gas entropy for $I_\phi$ and $I_\mathcal{K}$, a Boltzmann
configurational entropy for $I_\mathcal{V}$ and $I_\rho$, a connected mutual information for
$I_\mathcal{C}$. The coefficient that converts each to thermodynamic units is then the same
in every case,
\begin{equation}
  \alpha_\phi = \alpha_\rho = \alpha_\mathcal{V}
  = \alpha_\mathcal{K} = \alpha_\mathcal{C} = k_B,
\end{equation}
so the wake entropy collapses to
\begin{equation}
  S_{\text{wake}} = k_B\, W[\Psi] = k_B \sum_i I_i.
\end{equation}
The relative-weight problem flagged in \S\ref{sec:candidatefunctional} therefore dissolves:
once each channel is defined as a proper entropy, there are no free relative weights, and the
only dimensionful constant in the construction is $k_B$. The content did not disappear --- it
moved, from ``what are the relative weights'' to ``what is the correct entropy functional for
each channel,'' which is where it belongs. What remains genuinely open is finer: the
subleading terms in each functional, the regularization-dependent constants
($\mu_{\text{eff}}$, $v_0$, $g_{\text{int}}$), and the derivation of the channel-wise
coarse-grainings as a single renormalization flow rather than five separately specified maps.

\subsection{The Initial-Condition Question}

A monotonicity claim for $W[\Psi]$ has an immediate consequence: there must have been an
earliest moment at which $W$ was minimal. This is the SSV analogue of the past hypothesis in
standard cosmology — the assumption that the universe began in a low-entropy state, as
emphasized by Penrose's Weyl-curvature hypothesis~\cite{Penrose1979} and developed in
philosophical detail by Albert~\cite{Albert2000} and Price~\cite{Price1996}. The
present paper does not attempt to derive that condition. It treats it as a separate problem
deferred to a later cosmological paper in the SSV program, where the initial condition on
$\Psi$ should be addressed alongside the structure of the early medium. What the present
paper does claim is weaker and self-contained: \emph{given} a low-wake initial state of
$\Psi$, the dynamics described in \S4 and \S6 cannot locally undo wake accumulation, and
$W[\Psi(t)]$ therefore rises along physical histories. Whether the initial condition itself
is natural, fine-tuned, or derivable from deeper structure is a question outside the scope of
this paper.

\subsection{Two Quantities, Not One: Where the Observer Enters}\label{sec:twoquantities}

The preceding subsections actually describe two distinct quantities, and conflating them is
the source of the standard worry about whether entropy is ``in the world'' or ``in the head.''
The wake functional $W[\Psi]$ is a property of $\Psi$ alone: it is observer-independent, in
the same sense that the density $\rho(\mathbf{x},t)$ is observer-independent. The
coarse-grained entropy that follows from $\Omega_{\text{wake}}$ depends on the choice of
$(\ell,\mathcal{O})$, and is therefore observer-relative in the standard thermodynamic sense.

SSV takes both seriously. There is a structural burden carried by $\Psi$ regardless of
description, and there is a coarse-grained shadow of that burden read off by an observer who
resolves the medium only down to a particular scale and observable set. The two coincide in
the limit where $(\ell,\mathcal{O})$ resolves every wake channel; they diverge wherever they
do not. The monotonicity argument of \S\ref{sec:nottautological} belongs to the first
quantity. The recovery of ordinary thermodynamics in \S\ref{sec:extensivity} belongs to the
second.

\subsection{Possible Empirical Signatures}

The present paper does not advance a closed prediction, but the wake formulation suggests
where wake entropy could in principle diverge from standard thermodynamic entropy. The most
promising regimes are those where the medium's vortex and Kelvin-wave sectors carry a
disproportionate share of $W[\Psi]$ relative to thermal degrees of freedom.

Three illustrative regimes are worth naming. (i) Decaying vortex tangles in superfluid
$^4$He~\cite{Vinen2002}: in the late-time decay of quantum turbulence, conventional
thermodynamic entropy counts only the phonon bath, while $\alpha_{\mathcal{V}}
I_{\mathcal{V}} + \alpha_{\mathcal{K}} I_{\mathcal{K}}$ would track the residual tangle and
Kelvin-wave content that has not yet thermalized. A measured discrepancy between bulk thermal entropy and a
geometry-aware estimator over the same volume would be evidence for the wake decomposition.
(ii) BEC reconnection cascades~\cite{Villois2017}: the controlled reconnection of vortex
filaments in Bose-Einstein condensates allows direct imaging of the topological event and its
outgoing radiation, and is therefore a natural laboratory for comparing event-by-event wake
deposition against post-reconnection thermal balance. (iii) Two coarse-grainings of the same
state: any setting in which $(\ell, \mathcal{O})$ can be varied while the underlying $\Psi$
is held fixed should exhibit the predicted dependence of coarse-grained entropy on the
resolved wake channels, in the manner described in \S\ref{sec:boltzmann}.

None of these constitute a quantitative prediction in the present paper. They mark the
regimes where the formalism developed here would most plausibly first be tested.

\section{Formal Time Reversal and Physical Reversal}

A central distinction is needed between formal and physical reversal.

\subsection{Formal Reversal}

In many equations, one may perform the transformation
\begin{equation}
  t \to -t
\end{equation}
and obtain a formally related trajectory. This is a statement about the symmetry structure of
the equations.

\subsection{Physical Reversal}

Physical reversal would require the medium itself to perform the inverse history: gather the
distributed wake, reconstruct the exact incoming phases and correlations, and restore the prior
global state. In the language of \S3, it would require $\mathcal{R}[\Psi]$ to act with the sign
and content of an inverse update rather than a forward one — not merely a sign change on an
external time coordinate, but a physical operation the medium must actually carry out.

The existence of a formal reverse solution does not imply the existence of such an operation.
Without this distinction, the phrase ``time-reversal symmetry'' invites a misleading inference:
that because an equation admits a formally reversed trajectory, the physical world must be able
to run backward in the same sense. Formal reversal is a property of a mathematical description;
physical reversal is a claim about what operations the medium can actually perform.

\subsection{Boundary Data and the Inverse Problem}

The formal reverse of a dissipated event is not usually generated by a small local adjustment.
It requires exquisitely tuned incoming boundary data across the region into which the wake has
spread. In ordinary language, one would have to arrange the phonons, phases, correlations, and
vortex perturbations so that they converge on the event site in precisely the configuration
that reconstructs the former state.

SSV therefore locates irreversibility in the gap between a formal inverse trajectory and an
available physical operation. The reverse trajectory may exist in the mathematical solution
space. The present medium does not thereby contain a machine that prepares its boundary
conditions.

This is also where the standard Loschmidt~\cite{Loschmidt1876} and
Poincaré~\cite{Poincare1890} objections are absorbed. Loschmidt's paradox notes that
microreversible dynamics admit a formally reversed trajectory for every forward one; SSV does
not deny that, but holds (with Price~\cite{Price1996}) that the reverse trajectory is realized
only when the required incoming boundary data are actually supplied — which the medium does
not do spontaneously. Poincaré recurrence guarantees return to an arbitrarily close state in a
bounded system with discrete spectrum on cosmologically long timescales, but $\Psi$ in SSV is
effectively unbounded and supports continuous-spectrum excitations (outgoing phonons,
radiated Kelvin content, long-range correlations). The wake propagates into channels with no
finite recurrence time on physically relevant scales. Both objections are about what the
mathematics permits in the limit; the SSV claim is about what the medium delivers in finite
time.

\subsection{Classical Time Reversal Versus Quantum Unitarity}

The phrase ``time reversal'' covers two different mathematical structures, and the wake
argument is closer to one than the other. Classical time-reversal symmetry takes a trajectory
$\{q(t),p(t)\}$ to $\{q(-t),-p(-t)\}$ and is the symmetry Loschmidt invoked. Quantum unitary
evolution is more delicate: a closed system's state evolves under $U(t) = e^{-iHt/\hbar}$, and
$U^\dagger(t)$ is the formal inverse — backward evolution exists in principle for any pure
state.

The wake argument applies to both, but for related-yet-distinct reasons. Against classical
reversal it is the Loschmidt response just given. Against quantum reversal it is closer in
spirit to einselection and decoherence~\cite{Zurek2003,JoosZeh2003}: the wake of an event in
$\Psi$ is information entangled with environmental degrees of freedom across a vast region
of the medium, and the local subsystem no longer admits a pure-state description compatible
with the pre-event state. SSV does not deny global unitarity; it claims that the globally unitary operation
needed to physically reverse an event is not available to any local agent in $\Psi$, because
the requisite environmental boundary conditions cannot be locally prepared. The distinction
between formal and physical reversal is, in the quantum case, the distinction between
the existence of $U^\dagger$ as an operator and the existence of an apparatus that implements
it on a wake-dispersed state.

\subsection{Application: Closed Timelike Curves and Quantum CTCs}

The wake argument has direct consequences for physical proposals that invoke time travel as
a literal possibility. Two cases are worth naming.

\emph{Classical closed timelike curves.} General relativity admits solutions
(Gödel~\cite{Godel1949}, Kerr-interior, Tipler cylinder, traversable wormholes with exotic
matter~\cite{MorrisThorne1988}) containing closed timelike curves. The usual interpretation is that a worldline could in principle return to
its own past. In SSV this reading is incoherent for the reason developed in \S3: there is no
past region to return to. A closed timelike curve in the metric, if such a configuration of
$\Psi$ were physically realizable at all, would have to be reinterpreted as a closed structure
in the present state of the medium — a peculiar geometry, not an access route to an earlier
moment. The metric solution may exist; what does not exist is a destination at the other end
of the loop.

\emph{Deutsch's quantum CTCs.} Deutsch's D-CTC proposal~\cite{Deutsch1991} restores
self-consistency by demanding a fixed-point density matrix on the chronology-violating
region:
$\rho_{\text{CTC}} = \operatorname{Tr}_A[U(\rho_A \otimes \rho_{\text{CTC}})U^\dagger]$. The
mathematical content is a fixed point of a quantum channel; it is closed under any unitary
$U$ and famously implies computational consequences (efficient solution of \textsc{NP}-hard
problems, distinguishability of non-orthogonal states; see also the postselected variant of
Lloyd~\textit{et~al.}~\cite{Lloyd2011}). The SSV response is structural rather
than computational: the D-CTC prescription quietly assumes that the fixed-point state is
\emph{prepared} on the chronology-violating region, but preparation is itself a wake-writing
operation in $\Psi$. There is no agent in the medium whose preparing the fixed-point state
does not itself write distributed wake whose pre-state cannot be reassembled. Deutsch's
formalism solves the consistency problem at the level of operators while ignoring the
ontological problem at the level of $\Psi$: the wake of preparing the loop is already
elsewhere by the time the loop is supposed to close on itself.

Both cases share the same structure. They treat a formal solution (a CTC metric, a fixed-point
density matrix) as if the formal existence of the solution sufficed to make the physical
operation available. SSV holds that this is the same category error as before, dressed in
heavier mathematics.

\section{Discussion and Open Problems}

The paper makes a narrow ontological claim. The wake functional $W[\Psi]$ has been carried
from a schematic checklist to a complete channel set: all five channels are defined as
concrete dimensionless functionals, all coefficients are calibrated, and each channel reduces
to a recognizable entropy (\S\ref{sec:phononreduction}--\S\ref{sec:channelcomplete}). Several
technical tasks nonetheless remain open:

\begin{enumerate}
  \item pin down the subleading terms in each channel functional --- loop-closure and
  loop-number corrections in $I_\mathcal{V}$, interaction corrections in the soliton-gas
  $I_\rho$, the precise connected-part subtraction in $I_\mathcal{C}$ --- and compute the
  regularization-dependent constants $\mu_{\text{eff}}$, $v_0$, and $g_{\text{int}}$ from the
  microscopic field equation rather than leaving them as order-unity inputs,
  \item derive the five channel-wise coarse-grainings as a single renormalization flow on
  $\Psi$, rather than five separately specified $(\ell,\mathcal{O})$ maps,
  \item quantify the role of reconnection within the broader wake picture by measuring
  per-event wake deposition against the channel decomposition of \S\ref{sec:channelfate},
  \item connect this time framework to gravity and black holes in later papers without
  overloading the present one.
\end{enumerate}

The gain from the present framing is conceptual clarity. If time is local change, and the past
survives only as present wake, then irreversibility is not an added rule imposed on a reversible
block universe. It is the physical consequence of how events are written into a real medium.

\section{Conclusion}

The payoff of the present framing is that irreversibility stops being an extra rule layered on
top of a reversible block universe. If time is local change in $\Psi$ and the past survives
only as wake structure in the present state, then the asymmetry of experience and the
asymmetry of equations are no longer in tension: they describe different objects. Equations
describe trajectory classes that admit formal reversal; the medium carries out forward updates
whose distributed wake is not locally invertible. Entropy, on this account, is what that
non-invertibility looks like when it is read off the present state of $\Psi$.

What later papers inherit from this one is therefore a constraint, not a theorem. Any
SSV-derived account of gravitation, cosmology, or black-hole structure must be compatible with
the claim that the past is not a place — and any candidate entropy law for $\Psi$ must
ultimately reduce to a wake-functional of the kind sketched in \S5. The technical work of
turning that constraint into closed expressions is left to subsequent papers and to the open
problems listed above.

\begin{thebibliography}{99}

\bibitem{SSV-I}
S.~Norland, \textit{Saturated Superfluid Vacuum I: Topological Defects in a Saturated
Superfluid Vacuum --- The Geometric Origin of Matter and the $\alpha$-Harmonic Mass Ladder},
10.5281/zenodo.19651986, 2026.

\bibitem{SSV-II}
S.~Norland, \textit{Saturated Superfluid Vacuum II}, SSV program, 2026.

\bibitem{Barbour1999}
J.~Barbour, \textit{The End of Time: The Next Revolution in Physics}, Oxford University
Press, 1999.

\bibitem{Rovelli1993}
C.~Rovelli, ``Statistical mechanics of gravity and the thermodynamical origin of time,''
\textit{Class.~Quantum Grav.} \textbf{10}, 1549--1566 (1993).

\bibitem{Rovelli2018}
C.~Rovelli, \textit{The Order of Time}, Riverhead Books, 2018.

\bibitem{Price1996}
H.~Price, \textit{Time's Arrow and Archimedes' Point: New Directions for the Physics of
Time}, Oxford University Press, 1996.

\bibitem{Albert2000}
D.~Z.~Albert, \textit{Time and Chance}, Harvard University Press, 2000.

\bibitem{Penrose1979}
R.~Penrose, ``Singularities and time-asymmetry,'' in \textit{General Relativity: An Einstein
Centenary Survey}, eds.\ S.~W.~Hawking and W.~Israel, Cambridge University Press, 1979,
pp.~581--638.

\bibitem{Loschmidt1876}
J.~Loschmidt, ``Über den Zustand des Wärmegleichgewichtes eines Systems von Körpern mit
Rücksicht auf die Schwerkraft,'' \textit{Sitzungsber.~Akad.~Wiss.~Wien} \textbf{73},
128--142 (1876).

\bibitem{Poincare1890}
H.~Poincaré, ``Sur le problème des trois corps et les équations de la dynamique,''
\textit{Acta Math.} \textbf{13}, 1--270 (1890).

\bibitem{Zurek2003}
W.~H.~Zurek, ``Decoherence, einselection, and the quantum origins of the classical,''
\textit{Rev.~Mod.~Phys.} \textbf{75}, 715--775 (2003).

\bibitem{JoosZeh2003}
E.~Joos, H.~D.~Zeh, C.~Kiefer, D.~Giulini, J.~Kupsch, and I.-O.~Stamatescu,
\textit{Decoherence and the Appearance of a Classical World in Quantum Theory}, 2nd ed.,
Springer, 2003.

\bibitem{Godel1949}
K.~Gödel, ``An example of a new type of cosmological solutions of Einstein's field
equations of gravitation,'' \textit{Rev.~Mod.~Phys.} \textbf{21}, 447--450 (1949).

\bibitem{MorrisThorne1988}
M.~S.~Morris and K.~S.~Thorne, ``Wormholes in spacetime and their use for interstellar
travel: A tool for teaching general relativity,'' \textit{Am.~J.~Phys.} \textbf{56},
395--412 (1988).

\bibitem{Deutsch1991}
D.~Deutsch, ``Quantum mechanics near closed timelike lines,''
\textit{Phys.~Rev.~D} \textbf{44}, 3197--3217 (1991).

\bibitem{Lloyd2011}
S.~Lloyd, L.~Maccone, R.~Garcia-Patron, V.~Giovannetti, Y.~Shikano, S.~Pirandola,
L.~A.~Rozema, A.~Darabi, Y.~Soudagar, L.~K.~Shalm, and A.~M.~Steinberg, ``Closed timelike
curves via postselection: Theory and experimental test of consistency,''
\textit{Phys.~Rev.~Lett.} \textbf{106}, 040403 (2011).

\bibitem{Vinen2002}
W.~F.~Vinen and J.~J.~Niemela, ``Quantum turbulence,''
\textit{J.~Low Temp.~Phys.} \textbf{128}, 167--231 (2002).

\bibitem{Villois2017}
A.~Villois, D.~Proment, and G.~Krstulovic, ``Universal and nonuniversal aspects of vortex
reconnections in superfluids,'' \textit{Phys.~Rev.~Fluids} \textbf{2}, 044701 (2017).

\end{thebibliography}

\end{document}
